\lecture{5}{Tue 03 Feb 2026 13:58}{Concluding motivation with the divergence complex}

\begin{proof}
	Suppose we have vector spaces $M,N$, and subspaces $K\subseteq M$ and $L\subseteq N$. Recall that $F : M \to N$ descends to $M/K \to N/L$ whenever $F(K)\subseteq L$. We apply this to the case $P(C^\infty (U_1))\otimes P(C^\infty (U_2))\to P(C^\infty (V))$. The tensor product admits a projection $q_{12} $ to $\operatorname{Obs} ^q(U_1)\otimes \operatorname{Obs} ^q(U_2)$. We can then write
	\[
		\Ker q_{12} =(\Im (\operatorname{Div}_{U_1} )\otimes P(C^\infty (U_2))+(P(C^\infty (U_1))\otimes (\Im \operatorname{Div} _{U_{2} } ))
	.\]
	So we want to induce a map $\widetilde{\pi } $ in the diagram
	\[\begin{tikzcd}[row sep = 2.5em, column sep = 4em]
		\Ker q_{12} \ar[hook, r]  & P(C^\infty (U_1))\otimes P(C^\infty (U_2))\ar[" q_{12}  ", r] \ar[" \pi  ", d]  & \operatorname{Obs} ^q(U_1)\otimes \operatorname{Obs} ^q(U_2) \ar[dashed, " \pi  ", d]  \\
	\Im \operatorname{Div}_V \ar[hook, r]  & P(C^\infty (V))\ar[" q "', r]  & \operatorname{Obs} ^q(V)
	\end{tikzcd}\]
	It's enough to check that $\pi $ identifies the kernel into the image of $\Im \operatorname{Div} _V$; we may freely split the kernel at the $+$.

	There is a key identity. Let $X$ a vector field on $N$, $f$ a function on $N$, and $\omega $ a measure. Then $\operatorname{Div} _\omega (fX)-f \operatorname{Div} _\omega \left(X \right) =Xf$. To prove this (in finite dimensions), recall that $\operatorname{Div} _\omega X =\mathcal{L} _X\omega / \omega $. Then Cartan's magic formula and the product rule yield (since $\omega $ is a top form and $\mathrm{d} \omega =0$)
	\[
		\operatorname{Div} _\omega X=\frac{\mathrm{d} \iota_{fX} \omega }{\omega } -f \frac{\mathrm{d} \iota _X\omega }{\omega } = \frac{\mathrm{d} (f\iota _X\omega )}{\omega } -f\frac{\mathrm{d} \iota _X\omega }{\omega } 
	\]
	\[
		=\frac{\mathrm{d} f \cdot \iota _X\omega }{\omega } +\frac{f\cdot \mathrm{d} \iota _X\omega }{\omega } -f\frac{\mathrm{d} \iota _X\omega  }{\omega } =\frac{\mathrm{d} f \cdot  \iota _X\omega }{\omega } 
	.\]
	Passing to coordinates, let $\omega =g(x_i)\mathrm{d}x^i$ and $X=h_j \partial _j$. Then
	\[
		\iota _X\omega =\sum_{j=1}^{n} (-1)^jh_jg \mathrm{d} x^1 \wedge \cdots \wedge \widehat{\mathrm{d} x^j} \wedge \cdots \wedge \mathrm{d} x^n
	.\]
	Then
	\[
		\mathrm{d} f \wedge \iota _X\omega =\sum h_j g \frac{\partial f}{\partial x_j} \mathrm{d} ^nx
	.\]
	Then $\mathrm{d} f \cdot \iota _X\omega /\omega $ is simply $\sum_{j} h_j \partial _jf = X(f)$. This proof used finite dimensions, but the statement still follows for infinite dimensional manifolds. Hence if $X\in \operatorname{Vect} (C^\infty (U))$ and $F\in P(C^\infty (U))$, then $\operatorname{Div} (FX)-F \operatorname{Div} (X)=X(F)$.

	We will prove the statement in the dense subspaces $\widetilde{P} $ and $\widetilde{\operatorname{Vect} } $ since it is moer clear why this follows. Take $F=f_1 \cdots f_r$ and $X=g_1 \cdots g_s \frac{\partial }{\partial \phi } $, with $f_i\in C_c^\infty (U_1)$ and $g_i,\phi \in C^\infty _c(U_2)$. Then
	\[
		X(F)=\sum_{j=1}^{r} g_1 \cdots g_s f_1 \cdots \widehat{f} _j \cdots f_r \int_Vf_j\phi \,\mathrm{d} vol
	.\]
	But since $f_j,\phi $ have disjoint supports, the integral vanishes. Thus $X(F)=0$, showing that $\operatorname{Div} (FX)=F \operatorname{Div} (X)$. This implies that if something is a divergence then it all is???

	If divergence were an ideal in general, then $Xf$ would always be a divergence, which need not hold.
\end{proof}

This gives $\operatorname{Obs} ^q$ the structure of a precosheaf, together with maps
\[
	\left\{ U_i\subseteq V \right\} _{1\leq i\leq r} \text{ disjoint}:\qquad \bigotimes_{i=1}^r \operatorname{Obs} ^q(U_i)\to V
.\]
Hence $\operatorname{Obs} ^q$ is a prefactorization algebra.

In physics, there is a parameter $\hbar $ which determines the strength of quantum effects. When $\hbar =0$, the system is classical, but for nonzero $\hbar $ (as is reality) we see quantum effects. In finite dimensions (purely for illustration), we can take a Gaussian measure
\[
	e^{-S/\hbar } \mathrm{d} ^nx
.\]
Then
\[
	\operatorname{Div} _\omega \left(\sum_{i} f_i \partial _i\right)=-\frac{1}{\hbar } \sum_{i=1} ^n f_i \frac{\partial S}{\partial x_i} + \sum_{i=1}^{n} \frac{\partial f_i}{\partial x_i} 
.\]
Hence polynomials will simply be polynomials mod the image of this divergence, and when $\hbar \neq 0$, this is the same as the image of $\hbar \operatorname{Div} $. This allows us to clear the denominator and take $\hbar \to 0$:
\[
	\operatorname{Div} _\omega \left( \sum_{i} f_i \partial _i \right) =-\sum_{i=1}^{n} f_i \frac{\partial S}{\partial x_i} +\hbar \sum_{i=1}^{n} \frac{\partial f_i}{\partial x_i} \xlongrightarrow{\hbar \to 0} -\sum_{i=1}^{n} f_i \frac{\partial S}{\partial x_i} 
.\]
The set of all $-\sum_{i}f_i \partial_iS$ is the ideal of $P(\mathbb{R} ^n)$ generated by partials of $S$. Hence in classical mechanics in finite dimenisons, quantum observables go to
\[
	P(\mathbb{R} ^n) / \left\langle \partial _iS \right\rangle_{1\leq i\leq n} 
.\]
This is called the Jacobian ideal. Recall that although quantum mechanics is about path integrals (polynomials modulo divergences), classical mechanics was about critical loci. Note that by setting $\partial _iS=0$, we obtain only elements of the critical locus of $S$. Hence $P(\mathbb{R} ^n) / \left\langle \partial _iS \right\rangle $ is polynomial functions on the critical locus of $S$, meaning we recover $\mathcal{O} ^{cl} (\operatorname{Crit} (S))$. The same follows for infinite dimensional free field theories. I will not explain it because I'm tired of typing, but the various analogies of the infinite-dimensional setting to Gaussian measures should make the analogue clear.

This is a very simple example, and oftentimes what we will get is not very nice. Later on we will improve our construction by taking a differential graded object as a resolution of obserables, usually a chain complex, whose zeroth cohomology is observables. In the next chapter, we will properly define prefactorization algebras, and show that there is a family $\mathcal{O} _\hbar $ of flat prefactorization algebras over $\mathbb{C} [\hbar ]$, such that when $\hbar =1$ we get quantum observables, and in $h=0$ we get classical observables.

Our permanent friday room (except one weird day in March) is CDS 548. Hopefully we have 
